\section{Literature Review}

Memory-type charting began with the CUSUM of \citet{page_continuous_1954} and the geometric moving
average of \citet{roberts_control_1959}. For several correlated characteristics, two traditions
developed. The first extends univariate charts through Hotelling's $T^2$ statistic, culminating
for the memory-type case in the MEWMA chart, whose limit is set to achieve a specified $\ARL_0$ and
whose ARL depends on a mean shift only through its Mahalanobis distance
\citep{lowry_multivariate_1992}. The second, reviewed by \citet{macgregor_statistical_nodate},
projects the observations onto a few latent variables and monitors two complementary statistics:
$T^2$ on the retained scores and a squared prediction error (SPE) on the residual. Both traditions
rest on linear structure and on independence and normality assumptions whose failure can distort
run-length behaviour \citep{ahmed_statistical_2017}; distribution-free schemes relax the
parametric assumptions \citep{arslan_rankbased_2026}.

Representation learning offers a nonlinear counterpart to the projection step. Autoencoders learn
low-dimensional codes that outperform principal components for dimension reduction
\citep{hinton_reducing_2006}, convolutional autoencoders add locality and weight sharing
\citep{honkela_stacked_2011}, and one-dimensional convolutional networks suit sensor signals when
labelled data are scarce \citep{kiranyaz_1d_2021}. The analogy with latent-variable monitoring is
direct: the autoencoder plays the role of the principal component model and its reconstruction
error that of the SPE. For the latent space, one-class methods model normality when abnormal
examples are too scarce to model \citep{pimentel_review_2014}, and the OC-SVM scores observations
relative to an estimated support boundary \citep{scholkopf_estimating_2001}. Feeding learned
latent codes to an OC-SVM is established practice \citep{alaverdyan_regularized_2020,
pinon_one-class_2023,mabu_disaster_2019}, and \citet{guzman_sidewalk_2026} report that combining
autoencoder reconstruction error with an OC-SVM beats either component alone; all of these detectors,
however, classify each observation against a fixed threshold.

Hybrid ML and SPC research has followed two strands. In the first, a learner operates on charting
statistics while the chart remains classical: convolutional and recurrent networks recognise
control-chart patterns \citep{yu_control_2021}, deep networks and support vector regression model
EWMA-type statistics for censored reliability data \citep{lee_residual_2023,lee_monitoring_2024},
and models trained on classical charting statistics widen the range of detectable shifts
\citep{mukhtiar_enhanced_2025}. These learners need labelled out-of-control examples or a parametric
model from which training data are simulated. In the second strand, representation learning enters
the monitoring statistic itself: \citet{biegel_deep_2022} compute $T^2$ and SPE statistics from a
deep autoencoder's latent representation and reconstruction error on a real sheet-metal forming
process. That approach is unsupervised and acts on process observations, but its statistics are
memoryless and it is assessed by classification metrics rather than run lengths. Within the
literature reviewed, no study combines a CAE reconstruction error with an OC-SVM score on the CAE
representation in a memory-type chart whose limit is calibrated to a nominal $\ARL_0$ and compared
with a classical MEWMA under the same target; this study addresses that gap.
